\documentclass[11pt]{article}
\title{A newer LinuxPTP ingress timestamp need not carry a newer estimate}
\begin{document}
\maketitle

The pinned LinuxPTP management report \texttt{TIME\_STATUS\_NP}
copies the stored master offset and the sync ingress timestamp.
Those fields need not refer to one successful offset calculation.
An ingress timestamp alone is not a certificate of estimation age.

\paragraph{Executed selected code.}
The synchronization method stores the newly delivered ingress time
before asking \texttt{tsproc\_update\_offset} for an estimate.
In filter mode without a valid filtered path delay, that method
returns an error before writing to the offset pointer. For the
non-free-running case, the synchronization prefix then returns
without undoing the ingress write. The management switch case
copies both fields as they stand. This is separate from the
source's explicit \texttt{step\_window} rejection, which returns
before the ingress write.

The replay starts with stored offset 6 and ingress 100. A sync
with ingress 1120 reaches that filter path. Without a valid delay,
the report is now offset 6 with ingress 1120, even though no offset
was computed at that ingress. With valid delay 114 and the same
origin and ingress, the source recomputes offset 6 and reaches
the later synchronization code. At the compiled seam the
management fields are identical between the two branches.
A valid delay of 113 produces offset 7 as a positive control.
The report identity and other exported metadata are held fixed
in this fixture. No setter is replaced in the selected prefix
or export case. The surrounding servo, actual timers, wire
serialization and network are not part of the replay.

\paragraph{Implication for the risk audit.}
A consumer requiring an offset estimate at the acquisition cut
must establish that the timestamp and value belong to the same
successful calculation, or otherwise bound the retained estimate's
age. Even a fresh current measurement need not certify an entire
past acquisition epoch. A usable historical clock-class query
also needs the clock identity/epoch, a phase and rate envelope,
and error bounds backed by the actual clock and sensor sources.
The fixture does not calibrate those bounds or show a field
failure. It rules out using this report alone as a fresh historical
query in the stated code path.
\end{document}
